\subsubsection{The finite event profile}

The first event implementation supports consent and non-preemptive achievement
duties. Events have immutable identifiers, stream identifiers, occurrence and
recorded timestamps, and an authoritative sequence number. Timestamps are UTC
with microsecond precision; legal date conversion explicitly uses
\texttt{Asia/Hong\_Kong}. Equal occurrence times can use sequence order only
when the connector declares that ordering authority. Otherwise the result is
\texttt{E\_ORDER\_UNRESOLVED}. A duplicate event with identical bytes is a
no-op; reuse of its identifier for different content is an error.

Consent begins as not established only when the stream is complete from its
declared inception, or is complete after a reviewed initial snapshot. An incomplete export
instead yields an unknown history. A grant creates an active consent generation;
withdrawal makes it withdrawn; a new grant needs new evidence and creates a new
generation. Observed transactions remain in the audit history even when the
current consent state would prevent streamlined use.

An obligation records actor, action, subject, source version, activation and a
deadline or explicit absence of a deadline. Let $a$ be activation and $d$ the
exclusive deadline. Under this proposed profile, a matching performance at time
$p$ is timely when
\begin{equation}
 a\leq p<d.
 \label{eq:timely}
\end{equation}
This is the chosen boundary convention for the profile, not an assertion that
every SFC deadline has the same wording. A local due date is converted to the
next local midnight. An anniversary requires an explicit leap-day policy; a
reviewed choice of February 28, March 1 or rejection replaces an implicit library
default. Business-day calendars are deferred.

Absence of a timely record establishes breach only when a watermark shows that
the observation window through $d$ is complete. The wall clock cannot create
that evidence. Without completeness the obligation stays open. Timely performance
produces \texttt{SATISFIED\_ON\_TIME}; a complete window without it produces
\texttt{BREACHED}. A subsequent matching performance produces
\texttt{PERFORMED\_LATE} and retains the breach. A deadline-free request can be
performed but does not become overdue through an invented number of days.

For example, consider a synthetic review activated at midnight UTC on September
1 and due at midnight on September 2. A record one microsecond before the latter
instant is timely; a record exactly at it is late. The fixtures contain both.
These dates illustrate the runtime convention rather than an SFC-prescribed
review deadline. If a late-arriving record later proves timely performance,
replay produces a corrected assessment with a new knowledge cutoff; the original
assessment remains visible. Maintenance duties, pre-emptive satisfaction,
waivers and legal compensation require future profiles and are rejected now.
